\chapter{Introduction}
\label{sec:introduction}

\section{contribution}
\begin{itemize}
    \item I test both the FE-MSM and SeqGPLVM methods on a same 
    synthetic data and compare their performance. 
    \item by presenting these two methods in a same paper, 
    I can also highlight the differences and similarities between them 
    and perhaps provide some insights about when one method might be preferred over the other. and also 
    encourage future research to explore the connections between methods developed in different fields.
    \item I also make the assumptions behind these two methods more explicit and perhaps discuss 
    about the plausibility of these assumptions in real world settings. 
    \item I provide replicable code for both methods and make it available to the public. 
    \item connection with the blessing of multiple causes paper by Wang and Blei. specifically, 
    the analogous assumption of no single cause confounding in the time varying treatment setting 
    and also the critique made by D'amor about the deconfounder paper 
    and how it applies to the time varying treatment setting. 
    \item related to the previous point, in the paper towards clarifying the
    theory of deconfounder by wang and Blei, they accepted that identifying the potential 
    outcomes for the complete sets of treatments is not possible but it is 
    possible to identify the potential outcomes for subsets of treatments. 
    This is similar to the fe-msm approach where we only identify the potential outcomes for subsets of treatments 
    but they did not explicitly discuss this point in their paper. I can make this connection more explicit. 
\end{itemize}

    